\scanpage{53}
\mspage{12}
\unsure{relation déjà notée|résultat déjà noté|relations déjà notées} plus haut

\deleted{a) Cas de $r_0$ --- Il provient des relat\textsuperscript{ons} \unsure{$r$}}

\deleted{$s_{r_0}(\unsure{\zeta_{ij}|\xi'_j})=$}

Soit $W$ le groupe d'autom de $(\widehat E,\widehat\varphi,\lambda)$ \inserted{(\unsure{le même|le mien|ou ce qui revient au même} de $\Delta$)} engendré par les $s_r$. Je dis que $W$ \uline{opère transitivement sur l'ens des racines}, donc aussi sur l'ens des parties \inserted{$(\subset\mathfrak P_6(\Delta))$} de $\Delta$ associées --- qui sont donc toutes des parties basiques. Pour ceci, notons d'abord que
\[
\begin{aligned}
&\deleted{s_{r_{ij}}\,\zeta'_k=\left\{\zeta_k\quad k\ne i,j\right.}\\
&s_{r_{ij}}\text{ est l'autom de }\Delta_0\ \deleted{\text{engendré par}}\ \inserted{\text{associé à}}\ \text{l'élém}^{\rm t}\\
&\qquad\text{de }\unsure{\mathfrak S|\mathscr S}_6,\quad\deleted{\text{permutation}}\ \inserted{\deleted{\text{une}}\ \text{transposition}}\ \text{de }i\text{ et }j.
\end{aligned}
\]
(vérification immédiate) donc

\uline{Corollaire} $W$ contient $\mathfrak S_6$

\uline{Corollaire} les $r_{ij}$ sont conjugués par $W$, et les $r_{ijh}$ sont conjugués par $W$.

Il reste donc à voir que $r$ est conjugué à un $r_{ij}$ et à un $r_{ijh}$. \deleted{Car} Or on vérifie \deleted{\unsure{$r_\Delta$}} de suite
\[
s_{r_{ijh}}(r)=r_{\ell,m,n}
\]
donc $r$ conjugué aux $\unsure{r_{\ell mn}|r_{\ell,m,n}}$, et
\[
s_{r_{ijh}}(r_{i\unsure{\ell|g}})=r_{jh\ell}
\]
donc \deleted{\unsure{donc} $r_{i\ell}$ \unsure{est}} un des $r_{\ldotp\ldotp}$ ($r_{i\ell}$) conjugué à un des $r_{\ldotp\ldotp\ldotp}$ ($r_{jh\ell}$) et on a gagné.

On sait donc plus précisément, compte tenu que $W\supset\mathfrak S_6$

\uline{Corollaire} $W$ est transitif sur l'ens des \deleted{bases} racines, sur l'ens des bases \inserted{(qui sont en corr. $1$-$1$ avec les racines)}, sur l'ens des bases \uline{ordonnées}, et sur $V$ ens \deleted{des} des couples $(\delta',\delta)$ de deux sommets non liés. Enfin, $W$ est identique au groupe de tous les automorphismes de $\Delta$.
% archive sheet number: 52
